\documentclass[11pt,letterpaper]{article}
\usepackage[margin=.89in,headheight=15pt]{geometry}
\usepackage[T1]{fontenc}
\usepackage{lmodern,amsmath,amssymb,amsthm,mathtools,booktabs,enumitem,microtype,fancyhdr,xcolor,hyperref,array,longtable}
\hypersetup{colorlinks=true,linkcolor=[rgb]{.50,.31,.08},urlcolor=[rgb]{.50,.31,.08},pdftitle={GRID-K1: Infrastructure Continuity Geometry, Exact Cut Certificates and Admission Proofs},pdfauthor={Terry Allen Snyder}}
\definecolor{gold}{RGB}{125,84,31}\definecolor{muted}{RGB}{85,75,66}
\newtheorem{theorem}{Theorem}[section]
\newtheorem{lemma}[theorem]{Lemma}
\newtheorem{proposition}[theorem]{Proposition}
\newtheorem{corollary}[theorem]{Corollary}
\theoremstyle{definition}\newtheorem{definition}[theorem]{Definition}\newtheorem{example}[theorem]{Example}
\theoremstyle{remark}\newtheorem{remark}[theorem]{Remark}
\pagestyle{fancy}\fancyhf{}\lhead{\small\textcolor{muted}{V.A. / DEAGLE MATHEMATICS}}\rhead{\small\textcolor{muted}{Public research edition}}\cfoot{\small\thepage}
\setlength{\parindent}{0pt}\setlength{\parskip}{5pt}
\setlist[itemize]{topsep=3pt,itemsep=2pt,leftmargin=*}
\setlist[enumerate]{topsep=3pt,itemsep=3pt,leftmargin=*}
\newcommand{\R}{\mathbb R}\newcommand{\dd}{\mathrm d}\newcommand{\pos}{\mathbb R_{\ge0}}\newcommand{\ind}{\mathbf 1}
\begin{document}
\begin{center}
{\small\color{gold}\MakeUppercase{Deagle Mathematical Research / Formalization 002}}\\[10pt]
{\LARGE\bfseries GRID-K1: Infrastructure Continuity Geometry}\\[5pt]
{\large Admissible Transitions, Exact Cut Certificates,\; Cumulative Debt and Proof Receipts}\\[11pt]
{\large Terry Allen Snyder}\\[4pt]
{\small V.A. / Deagle Mathematical Research \quad 11 November 2026}\\[8pt]
{\small\color{muted}Source-preserving mathematical derivation | public research edition | v1.0}
\end{center}
\vspace{5pt}\hrule\vspace{11pt}
\begin{abstract}
The author-held \emph{World Kernel Atlas} identifies GRID-K1 as an infrastructure continuity kernel, and the separate \emph{Critical Infrastructure Continuity Kernel v1.0} describes authority-bound continuation, degradation, rerouting, escalation, halt and refusal across coupled power, water, telecommunications and civic services. This manuscript extracts a precise, testable mathematical layer without treating the source design proposal as already-proved deployment. We introduce a finite directed service-capacity network with dependency constraints, nonnegative cumulative debt, freshness and authority predicates, and a pre-effect admissibility guard. We prove a necessary cut-capacity obstruction, give a constructive sufficiency certificate through admissible flows, prove conditional invariant preservation under enforced guards, and show no return to a strict debt reserve absent an explicit reset law. A frozen seven-case synthetic example is computed with exact rational arithmetic, seven proposition-scoped certificates and seven chained receipts. These results certify the specified finite model and example; they do not establish the behavior of a physical electrical grid, an installed protection controller or an unrestricted universal cascade-prevention theorem.
\end{abstract}
\textbf{Source status and custody.} The parent infrastructure kernel is labeled \texttt{approval\_draft} in its original source. Both original Library documents are retained unchanged in the \emph{author-review} evidence archive. The \emph{public-safe} companion excludes those full drafts pending author disclosure approval; it includes their hashes and an itemized formula-to-source crosswalk. This is not a claim that the protected URK generator has been reproduced.

\section{Source problem and mathematical target}
The continuity proposal asks whether a coupled transition may take effect under live authority and interdependent state, and whether degradation or rerouting could shift a shortfall into an adjacent corridor. It identifies six decision categories, \textsc{Continue}, \textsc{Degrade}, \textsc{Reroute}, \textsc{Escalate}, \textsc{Halt}, and \textsc{Refuse}; only the first three represent positive service transitions in the finite model below. The remaining three are control dispositions whose legal and physical implementation are not implied by flow arithmetic.

We formalize one bounded question: \emph{Given a prepared graph, capacities, service demands, debt increments, dependency flags and trustworthy authority/evidence predicates, can a proposed mode be admitted without violating any of those declared mathematical constraints?} The example models generic transferable service capacity, not Kirchhoff power flow, voltage, frequency, inertia, protection coordination or restoration dynamics.

\section{Finite network and admissible continuity}
\begin{definition}[Coupled corridor network]\label{def:network}
Let $G=(V,E)$ be a finite directed graph with a distinguished source $s$, sink set $L\subset V\setminus\{s\}$ and capacities $c_e(m,\xi)\in\pos$ determined by a proposed mode $m$ and an attested state $\xi$. Let $q_\ell\ge0$ be the exact target service delivered to $\ell\in L$, and $f_e\ge0$ the proposed edge flow. Write $\delta^+(v)$ and $\delta^-(v)$ for outward and inward edges. A \emph{flow witness} satisfies
\begin{align}
0\le f_e &\le c_e(m,\xi) && (e\in E),\label{eq:cap}\\
\sum_{e\in\delta^-(v)}f_e-\sum_{e\in\delta^+(v)}f_e &= q_v &&(v\ne s),\label{eq:balance}\\
\sum_{e\in\delta^+(s)}f_e-\sum_{e\in\delta^-(s)}f_e &=\sum_{\ell\in L}q_\ell,\label{eq:source}
\end{align}
where $q_v=0$ for nonsinks. This abstraction assumes a single conserved nonnegative commodity, with no unmodeled store, loss, gain or incompatible commodity conversion.
\end{definition}
\begin{definition}[Cross-domain prerequisites and debt]\label{def:dependency}
Let $P\subseteq V\times V$ be a declared prerequisite relation: $(u,v)\in P$ means an active service at $v$ requires the attested service state of $u$ to be active. Let $z_v\in\{0,1\}$ denote the verified availability flag. For each monitored resource $e\in I\subseteq E$, let $d_e\ge0$ denote cumulative debt, $\Delta_e\ge0$ a declared increment and $B_e>0$ its budget. The prepared candidate must satisfy
\begin{equation}\label{eq:depdebt}
q_v>0\implies z_u=1\quad ((u,v)\in P),\qquad d_e^+=d_e+\Delta_e < B_e \quad (e\in I).
\end{equation}
The strict inequality is an operating \emph{reserve} convention; equality reaches the mathematical reserve boundary. It does not automatically constitute physical failure.
\end{definition}
\begin{definition}[Authority-bound admission guard]\label{def:guard}
Let $A(m,\xi)$ mean that a trusted authority checker has validated scope, signature, action rights and time window; let $F(\xi)$ mean that relevant evidence, synchronization and interlock attestations are sufficiently fresh. Let $M(m,q)$ denote prepared mode obligations: for \textsc{Continue} and \textsc{Reroute}, $q=q^{\rm base}$; for \textsc{Degrade}, $q\ge q^{\min}$ componentwise. Let $\mathcal C_{\rm flow}$, $\mathcal C_{\rm prereq}$ and $\mathcal C_{\rm debt}$ denote satisfaction of all flow constraints \eqref{eq:cap}--\eqref{eq:source} and the corresponding two conditions in \eqref{eq:depdebt}. Define
\begin{equation}\label{eq:guard}
\mathsf{Admit}(m,\xi,f,q,d^+)\iff
A(m,\xi)\wedge F(\xi)\wedge M(m,q)\wedge \mathcal C_{\rm flow}\wedge \mathcal C_{\rm prereq}\wedge \mathcal C_{\rm debt}.
\end{equation}
Every term is a stated Boolean condition, not a reference number standing in for a theorem. \textsc{Refuse}, \textsc{Escalate} and \textsc{Halt} are non-admission outcomes, not additional feasible flows.
\end{definition}
\begin{remark}[Limits of an attestation input]
A Boolean claim of ``fresh'' or ``authorized'' is not itself a cryptographic verification. Admission may rely on such a predicate only if a separately trusted implementation actually establishes it. Likewise, asserting a single governed path does not prove there is no bypass path in deployed equipment.
\end{remark}

\section{Cut obstruction, constructive feasibility and non-export}
\begin{theorem}[Necessary service-cut inequality]\label{th:cut}
Assume \eqref{eq:cap}--\eqref{eq:source}. For every subset $U\subset V$ with $s\in U$ and $U\cap L=\varnothing$,
\begin{equation}\label{eq:cut}
\sum_{\ell\in L}q_\ell\le\sum_{e\in\delta^+(U)}c_e(m,\xi).
\end{equation}
A violated cut is therefore an exact certificate that the specified full-demand flow cannot exist.
\end{theorem}
\begin{proof}
Sum the balance equations over $V\setminus U$. Internal edge flows cancel, leaving total demand equal to flow entering $V\setminus U$ from $U$ minus flow returning to $U$. Because return flows are nonnegative, and forward flows do not exceed their capacities, the total demand cannot exceed the outward cut capacity.
\end{proof}
\begin{corollary}[Full cut characterization in the stated single-commodity model]\label{cor:maxflow}
Add a super-sink $t$ and edges $\ell\to t$ with capacity $q_\ell$. If every $s$--$t$ cut in this augmented graph has capacity at least $Q=\sum_{\ell\in L}q_\ell$, then a flow satisfying all demands exists. For rational capacities a rational feasible flow exists.
\end{corollary}
\begin{proof}
By the max-flow/min-cut theorem, a flow of value $Q$ exists. The total capacity of the sink arcs equals $Q$, so each sink arc is saturated, delivering exactly $q_\ell$. With rational capacities, scaling by a common denominator reduces to an integral-capacity network, giving a rational flow after rescaling. This does not imply feasibility of multi-commodity or AC electrical-network constraints.
\end{proof}
\begin{proposition}[Local reroute checks are insufficient]\label{prop:export}
Checking capacity only on the \emph{new} path does not ensure a lawful reroute in a shared network. If all participating corridors are represented in $G$, however, a witness satisfying \eqref{eq:cap}--\eqref{eq:source} cannot overload any \emph{represented} edge or conjure net flow at a nonsource vertex.
\end{proposition}
\begin{proof}
The second assertion is exactly the componentwise capacity and node-balance equations. For the first, take a shared edge of capacity $5$ carrying $3$ existing units and a proposed rerouted $3$ units. The new route alone uses $3\le5$, but the combined requirement is $6>5$. Thus an isolated local check would accept a globally impossible transfer.
\end{proof}

\section{Invariant preservation and debt-bound continuity}
\begin{theorem}[Conditional pre-effect preservation]\label{th:invariant}
Let $X$ be the set of model states satisfying the prepared debt, dependency and service invariants. Suppose for each proposed transition $x\mapsto x^+$ (i) the admission guard checks the \emph{entire} next-state witness against $X$, (ii) the attested data are correct for that transition, and (iii) any permitted implementation follows the certified transition without an alternate actuation path. Then
\begin{equation}\label{eq:induction}
x_0\in X\quad\Longrightarrow\quad x_n\in X\quad\text{for every completed admitted step }n.
\end{equation}
\end{theorem}
\begin{proof}
At $n=0$ the hypothesis holds. If $x_n\in X$ and the next step is admitted, the guard has verified the asserted next state $x_{n+1}\in X$; the exact-execution hypothesis makes the executed next state equal to this witnessed state. Induction completes the proof.
\end{proof}
\begin{remark}[The critical deployment boundary]
Theorem \ref{th:invariant} is conditional, not an assertion that physical cascading failures are impossible. State estimation, model completeness, adversarial inputs, software atomicity, fail-safe implementation, protection equipment and unmodeled physical dynamics remain independent proof obligations. A negative gate decision also does not, by itself, establish that an immediate physical shutdown is safe.
\end{remark}
\begin{theorem}[No return of monotone debt to a strict reserve]\label{th:debt}
Let $d(t)=d(0)+\int_0^t k(u)\,\dd u$ with $k\ge0$ locally integrable and strict reserve $d<B$. If $\tau_B=\inf\{t:d(t)\ge B\}<\infty$, then $d(t)\ge B$ for every $t\ge\tau_B$ in the model interval. For discrete steps with $\Delta_e\ge0$, if $d_{e,n}\ge B_e$ at a step, no later nonnegative increment can restore $d_{e,n}<B_e$.
\end{theorem}
\begin{proof}
The integral, respectively each partial sum of nonnegative increments, is nondecreasing. Consequently debt cannot decrease below an already reached threshold. A service reset or repair mechanism would require an additional explicit state transition law and cannot be inferred here.
\end{proof}
\begin{corollary}[Debt reserve test for a prepared step]\label{cor:reserve}
If $d_e=3/4$, $B_e=1$ and $\Delta_e=1/8$, the next debt is $7/8<1$ and this particular guard passes. If $\Delta_e=3/8$, the next debt is $9/8>1$ and it fails. These are exact rational comparisons, not predicted equipment deterioration rates.
\end{corollary}
\begin{proof}Compute the sums in $\mathbb Q$.\end{proof}

\section{One frozen synthetic network: seven exact cases}
Consider the finite network shown by the edge-capacity table. The source $S$ routes a generic capacity resource through primary $A$, secondary $B$ and a normally disabled standby route $C$ toward water-service node $W$ and telecom-service node $T$. The labels identify coupled \emph{service obligations}; a unit is an abstract conserved unit, not a MW, water volume or bandwidth conversion. The seven-case contract instantiates no nonempty prerequisite relation $P$; it tests shared-network conservation, capacity, mode requirements, cumulative debt and declared authority/freshness inputs, not cross-domain prerequisite attestation.
\begin{center}\small
\begin{tabular}{@{}lrrrrrrrr@{}}\toprule
Edge & $S\!\to\!A$&$S\!\to\!B$&$S\!\to\!C$&$A\!\to\!W$&$A\!\to\!T$&$B\!\to\!W$&$B\!\to\!T$&$C\!\to\!T$\\\midrule
Normal cap.&6&5&0&4&2&4&3&0\\
Fault, standby off&3&5&0&4&2&4&3&0\\
Fault, standby on&3&5&4&4&2&4&3&4\\\bottomrule
\end{tabular}
\end{center}
The baseline obligations are $(q_W,q_T)=(6,5)$; the authorized degradation minimum is $(6,2)$. For the monitored $S\to A$ edge, $d_0=3/4$ with strict cap $B=1$. The proposed increments are $1/8$ except in the dedicated overbudget negative case.
\begin{example}[Baseline full service]
Use $f_{SA}=6$, $f_{SB}=5$, $f_{AW}=4$, $f_{AT}=2$, $f_{BW}=2$, $f_{BT}=3$. Each edge respects capacity and delivers $(6,5)$. Subject to fresh authority and the reserve check, \textsc{Continue} is admitted.
\end{example}
\begin{example}[Fault cut certificate]
After the primary capacity drops from $6$ to $3$ with standby disabled, the cut $U=\{S\}$ has outward capacity $3+5+0=8$, while baseline service requires $6+5=11$. By Theorem \ref{th:cut}, full-demand continuation is impossible; the deficit is exactly $3$ units. An unmodified $6$-unit flow along $S\to A$ also violates the primary edge capacity. The proposed action is \textsc{Refuse}.
\end{example}
\begin{example}[Authorized degradation]
With standby off and the fault present, set $f_{SA}=3$, $f_{SB}=5$, $f_{AW}=3$, $f_{BW}=3$, $f_{BT}=2$, and all other flows zero. This exactly delivers $(6,2)$ and respects the cut of capacity $8$. Provided the prepared degraded minima and authority are valid, the result is \textsc{Degrade}.
\end{example}
\begin{example}[Authorized reroute]
Activate the \emph{already modeled} standby $S\to C\to T$ corridor; set $f_{SA}=3$, $f_{SB}=5$, $f_{SC}=3$, $f_{AW}=3$, $f_{BW}=3$, $f_{BT}=2$, $f_{CT}=3$. The sinks receive $(6,5)$ and the source sends $11$ units through an available $12$-unit source cut. The entire shared graph is checked: \textsc{Reroute} is admissible if its activation right and live attestations hold.
\end{example}
Three further checks reject the same nominal reroute on stale evidence, reject a permitted degradation when the exact debt increment becomes $3/8$, and reject the full reroute when its activating permission is absent. Importantly, a mathematically feasible flow is \emph{not} automatically an admissible authorized action.

\begin{center}\small
\begin{tabular}{@{}p{.31\linewidth}p{.18\linewidth}p{.41\linewidth}@{}}\toprule
Prepared case & Decision & Precise basis\\\midrule
C01 Healthy baseline & \textsc{Continue} & 11 delivered, 11 source cut, fresh authority\\
C02 Fault/full demand & \textsc{Refuse} & Source cut 8 $<$ required 11; primary overload\\
C03 Fault/degraded & \textsc{Degrade} & 8 delivered, all edge and minima checks pass\\
C04 Fault/backup & \textsc{Reroute} & 11 delivered, source cut 12, full witness passes\\
C05 Stale backup & \textsc{Refuse} & Evidence freshness predicate false\\
C06 Debt exhaustion & \textsc{Refuse} & Next debt $9/8$ exceeds strict cap $1$\\
C07 Permission absent & \textsc{Refuse} & Reroute authorization not present\\\bottomrule
\end{tabular}
\end{center}
The exact result objects, contract hashes and receipt chain are frozen in the companion files. No case actuates or contacts live equipment.

\section{Prepared generation, certificates and reproducible replay}
The sequence is \emph{source custody $\to$ frozen mathematical contract $\to$ one identified exact run $\to$ seven arithmetic witnesses $\to$ seven certificates $\to$ seven chained receipts $\to$ separate read-only verification}. The public code is a transparent arithmetic reference implementation. It is not the protected URK parent generator or a production GRID-K1 runtime.

Each witness records the proposed flow, modeled capacity and conservation checks, authority/freshness predicates, mode minima, exact debt comparison and a source-cut certificate. Its certificate is deliberately limited to the proposition: \emph{the decision follows from the declared finite scenario and guard}. The receipt binds a run ID, certificate hash, prior receipt hash, recorded decision and failed constraints. A separate read-only program recomputes the rational constraints without importing the execution program, checks the source and contract digests when originals are present, validates the linked receipts and rejects a deliberately altered record. Hashes bind the frozen record against changes relative to its published digest; without an external trust anchor or signatures, replacing an entire bundle and all its hashes is not ruled out.

The source's words ``failure classes closed'' and ``no side path'' are design intentions until supported by separate installation, threat-model and physical integration tests. Seven example certificates do not establish those broader statements.

\section{Research claims, limits and examination invitation}
The contributions established here are: an explicit mathematical admission contract for a single-commodity coupled-service abstraction; a proof that a violated cut rules out a demanded flow; constructive exact rational flow witnesses for degraded and rerouted modes; and a conditional stepwise invariant theorem with cumulative-debt refusal. These are formalizations applying standard conservation, max-flow/min-cut and induction arguments, not claims that those classical theorems were newly invented.

\textbf{Not established:} a complete grid-state observer; AC/DC electric feasibility; operational security; signature verification; an installed no-bypass control path; field-safe halting; actual physical cascade suppression; universal sufficiency of this capacity abstraction; or an external certification. Physical grid validation requires its own equations, measurements and deployment proof.

A reviewer can challenge any exact step by citing the stated theorem, graph edge, scenario, contract field or receipt. The review request is for \emph{mathematical examination}, not for prior permission to continue research.

\appendix
\section{Original source-to-proof crosswalk}
\small
\begin{longtable}{@{}p{.24\linewidth}p{.37\linewidth}p{.27\linewidth}@{}}\toprule
Source object / assertion & Current mathematical carrier & Present standing\\\midrule\endhead
World Kernel Atlas: GRID-K1 & Infrastructure continuity across power, water, telecom; degrading service without unmodeled cross-domain harm & Original index declaration\\
Corridor profile & Graph, demands, dependencies, capacities, modes & Formal definition and prepared example\\
Authority envelope & $A(m,\xi)$ with permission, scope and validity & Assumed input; not cryptographic proof\\
Live-state surface & Capacity/mode flags and freshness predicate $F(\xi)$ & Modeled input, not field attestation\\
Maintenance debt & Monotone $d_e$ and strict reserve $d_e^+<B_e$ & Conditional theorem and rational examples\\
Continue / Degrade / Reroute & Positive admitted modes and conservation witnesses & Seven-case demonstration\\
Escalate / Halt / Refuse & Non-admission dispositions & \textsc{Refuse} modeled; other behaviors untested\\
Continuity receipt & Run/case IDs, source/contract hashes, linked certificates and receipts & Generated and replay-checked\\
``No side path'' / ``failure classes closed'' & Global enforcement of guard on every physical actuation path & Open implementation obligation\\
``Before cascading failure'' & Cut obstruction and conditional invariant preservation & Formal model only, not physical claim\\\bottomrule
\end{longtable}
\normalsize
\section{Reproduce the synthetic evidence}
With the included public-safe evidence folder:
\begin{verbatim}
python reproducibility/verify_exact_output.py
\end{verbatim}
This verifies the released deterministic arithmetic records and performs a tamper negative control. The author-only archive additionally contains the two original Library source files for source-byte verification. To intentionally produce another run, prepare a new contract with a new run ID and retain both versions; do not overwrite the frozen release.

\section{References and source boundary}
\begin{enumerate}
\item T. A. Snyder, \emph{Critical Infrastructure Continuity Kernel v1.0}, source-held JSON, Library record (2026), with status \texttt{approval\_draft}. Original SHA-256 in the manifest.
\item T. A. Snyder, \emph{World Kernel Atlas}, source-held DOCX, \texttt{GRID-K1} entry, Library record (2026). Original SHA-256 in the manifest.
\item T. A. Snyder, \emph{Admissibility--Debt Geometry: Boundary Hitting, Cumulative Debt, Irreversibility, and Exact Proof Witnesses}, Deagle public mathematical record (11 November 2026).
\item L. R. Ford, Jr., and D. R. Fulkerson, ``Maximal Flow Through a Network,'' \emph{Canadian Journal of Mathematics} \textbf{8} (1956), 399--404.
\end{enumerate}
\end{document}
